\chapter{ISTANTANEE DELL'ECOPOIESI}\label{istantanee-dell-ecopoiesi}

\hfill\break

Jason suggerì di prenotare le stanze a Cocoa Beach e di aspettarlo li
dove ci avrebbe raggiunto il giorno dopo. Stava facendo un ultimo giro
di domande e risposte con i media alla Perielio ma aveva cancellato
tutti gli altri impegni in programma prima dei lanci. Li voleva vedere
di persona senza che ci fosse una troupe della CNN a importunarlo con
domande stupide.

«Ottimo,» fece Diane quando glielo dissi, «così potrò fargli io tutte le
domande stupide che voglio.»

Ero riuscito a calmare le sue paure sullo stato di salute di Jason; no,
non stava morendo, e tutti quei buchi temporali sulla sua cartella
clinica erano affari suoi. Lo aveva accettato, o così sembrava, ma
voleva comunque vederlo, anche solo per tranquillizzarsi, come se la
morte di mia madre avesse scosso la sua fiducia nelle stelle fisse
dell'universo dei Lawton.

Così, usai il mio tesserino identificativo della Perielio e il mio
aggancio con Jase per affittare due stanze adiacenti in un Holiday Inn
con vista su Canaveral. Non molto tempo dopo l'ideazione del progetto
Marte, una volta annotate e ignorate le obiezioni dell'Agenzia per la
protezione dell'ambiente, una decina di piattaforme di lancio da acque
basse furono costruite e ancorate fuori dalla costa di Merritt Island.
Erano queste le strutture che potevamo vedere più chiaramente
dall'albergo. Il resto del panorama era fatto di parcheggi, spiagge
d'inverno e acqua azzurra.

Eravamo sul terrazzino della stanza di Diane. Aveva fatto la doccia e si
era cambiata dopo il viaggio in macchina da Orlando e stavamo per
scendere al piano di sotto e affrontare il ristorante della hall. Tutti
gli altri terrazzi che riuscivamo a vedere erano pieni di macchine
fotografiche e obiettivi: l'Holiday Inn era l'albergo destinato ai media
- Simon sarà stato pure diffidente verso gli organi di stampa laici ma
Diane all'improvviso c'era dentro fino al collo. Non riuscimmo a vedere
il Sole che tramontava ma la sua luce catturava le rampe e i razzi in
lontananza, facendoli sembrare più eterei che reali, uno squadrone di
robot giganti in marcia verso una qualche battaglia nella dorsale
medioatlantica. Diane stava a distanza dal corrimano del terrazzo come
se trovasse spaventosa quella vista.

«Perché ce ne sono così tanti?»

«Ecopoiesi a raffica» dissi.

Si mise a ridere, con una leggera aria di rimprovero. «È forse una delle
parole di Jason?»

Non lo era, non completamente. ``Ecopoiesi'' era una parola coniata da
un uomo di nome Robert Haynes nel 1990, ai tempi in cui la
terraformazione era solamente una speculazione scientifica. Tecnicamente
indicava la creazione di una biosfera anaerobica autoregolante dove non
ne era mai esistita nessuna, ma nell'uso moderno si riferiva a una
qualsiasi forma di modificazione puramente biologica di Marte.
L'inverdimento di Marte richiedeva due diversi tipi d'ingegneria
planetaria: la terraformazione rudimentale, ovvero l'innalzamento della
temperatura della superficie e della pressione atmosferica fino a una
soglia che rendesse possibile la vita, e l'ecopoiesi, in altre parole
l'utilizzo della vita microbica e vegetale al fine di bonificare il
suolo e ossigenare l'aria.

Lo Spin aveva già fatto il lavoro più pesante al posto nostro. Ogni
pianeta del sistema solare - tranne la Terra - era stato
significativamente riscaldato dal Sole in espansione. Ciò che rimaneva
da fare era il lavoro di fino dell'ecopoiesi. Ma esistevano molte vie
praticabili verso l'ecopoiesi, molti organismi candidati, dai batteri
che dimorano nelle rocce ai muschi alpini.

«Quindi si chiama raffica,» evinse Diane, «perché li manderete tutti.»

«Tutti, e quanti più possiamo permettercene, perché neanche su uno solo
di questi organismi c'è garanzia che si adatti e sopravviva. Però almeno
uno potrebbe farcela.»

«Anche più di uno.»

«E questo sarebbe un bene. Vogliamo un ecosistema, non una monocoltura.»

Infatti, i lanci sarebbero stati sincronizzati e scaglionati. La prima
ondata avrebbe portato solamente gli organismi anaerobici e
fotoautotrofi, forme di vita semplici che non avevano bisogno di
ossigeno e ricavavano energia dalla luce solare. Se fossero prosperati e
morti in quantità sufficiente avrebbero creato uno strato di biomassa
utile a nutrire ecosistemi più complessi. L'ondata successiva, a
distanza di un anno da allora, avrebbe introdotto organismi ossigenanti;
gli ultimi lanci senza equipaggio avrebbero compreso piante primitive
per preparare il suolo e regolare l'evaporazione e il ciclo delle
piogge.

«Sembra tutto così improbabile.»

«\emph{Viviamo} in tempi improbabili. Comunque no, non è detto che
funzioni.»

«E se non funziona?»

Scrollai le spalle. «Che abbiamo da perdere?»

«Un sacco di soldi. Un sacco di risorse umane.»

«Non mi viene in mente un uso migliore. Certo, questa è una scommessa, e
no, non è una cosa sicura, ma il risultato potenziale vale molto di più
del rischio. Ed è stato un bene per tutti, almeno finora. Un bene per il
morale della nazione e un buon modo di promuovere la cooperazione
internazionale.»

«Ma avrete ingannato tanta gente, molte persone comuni. Le avrete
convinte che lo Spin è una cosa che siamo in grado di gestire, qualcosa
che possiamo aggiustare con la tecnologia.»

«Intendi dire che abbiamo dato loro speranza.»

«Speranza del tipo sbagliato. E se fallite le lasciate senza alcuna
speranza.»

«Cosa vorresti che facessimo, Diane? Ritirarci sui tappetini da
preghiera?»

«Non sarebbe per nulla un'ammissione di sconfitta - la preghiera,
intendo. E se riuscirete, il passo successivo sarà mandare delle
persone?»

«Sì. Se riusciremo a rendere verde il pianeta, manderemo delle persone.»
Una proposta ancora più difficile e complessa dal punto di vista etico.
Avremmo mandato i candidati in gruppi di dieci. Avrebbero dovuto
resistere per un periodo imprevedibilmente lungo in alloggi
piccolissimi, con provviste limitate. Avrebbero dovuto subire,
all'impatto con l'atmosfera, una variazione di velocità impulsiva quasi
letale dopo mesi di assenza di gravità, seguita da una pericolosa fase
di atterraggio sulla superficie del pianeta. Se tutto ciò avesse
funzionato e se la loro scarsa dotazione di equipaggiamento di
sopravvivenza fosse atterrata al contempo e abbastanza vicino, avrebbero
dovuto acquisire da soli le competenze per sostentarsi in un ambiente
adatto all'insediamento umano solo in parte. Lo scopo della loro
missione non consisteva nel ritornare sulla Terra ma nel vivere
abbastanza a lungo da riprodursi in numero sufficiente e tramandare alla
loro progenie un modo di sopravvivere sostenibile.

«Quale persona sana di mente accetterebbe mai di farlo?»

«Ne rimarresti sorpresa.» Non potevo parlare per i cinesi, i russi o per
nessuno degli altri volontari internazionali, ma i candidati
nordamericani erano un gruppo di uomini e donne comunissimi. Erano stati
selezionati in base alla giovane età, alla costituzione robusta e alla
capacità di sopportare e tollerare ogni disagio. Solo alcuni erano stati
piloti collaudatori dell'Air Force ma tutti avevano quella che Jason
definiva ``la mentalità del pilota collaudatore'', una propensione
all'accettazione di un grave rischio fisico in nome di una conquista
eccezionale. E ovviamente, con ogni probabilità, erano quasi tutti
condannati, proprio come lo erano quasi tutti i batteri caricati su quei
razzi lontani. Il migliore risultato che potevamo ragionevolmente
aspettarci era che almeno un gruppo di sopravvissuti nomadi, tra quelli
che vagavano per i canyon muscosi delle Valles Marineris, potesse
incontrare un gruppo simile di russi, danesi o canadesi e generare
un'umanità marziana autosufficiente.

«E tu \emph{approvi} tutto questo?»

«Nessuno ha chiesto la mia opinione. Ma auguro loro il meglio.»

Diane mi lanciò uno sguardo come se volesse dire ``non mi basta'' ma
decise di non proseguire la discussione. Scendemmo nella hall con
l'ascensore per raggiungere il ristorante. Mentre ci mettevamo in fila
per il servizio al tavolo dietro una decina di tecnici dei notiziari
televisivi, probabilmente percepì l'entusiasmo crescente.

Dopo aver ordinato girò la testa per ascoltare i frammenti delle varie
conversazioni - parole come ``fotolisi'' e ``criptoendolitico'' e, sì,
``ecopoiesi'' - che traboccavano dai tavoli affollati, con i giornalisti
che provavano il gergo da usare per il giorno dopo o si sforzavano
almeno di comprenderlo. Si sentivano anche le risate e lo sbattere
incurante delle posate, c'era un'atmosfera di attesa vertiginosa per
quanto incerta. Era la prima volta dall'allunaggio di oltre sessant'anni
prima, che l'attenzione del mondo era stata focalizzata così a pieno su
un'avventura spaziale, e lo Spin le attribuiva quello che persino alla
spedizione lunare era mancato: un'urgenza reale e un senso di rischio
globale.

«Tutto questo è opera di Jason, vero?»

«Senza Jason e E.D. sarebbe comunque potuto succedere. Ma in modo
diverso, magari meno rapido ed efficiente. Jase è sempre stato al centro
del progetto.»

«E noi alla periferia. A orbitare attorno al suo genio. Ti svelo un
segreto. Ho un po' paura di lui. Ho paura di incontrarlo dopo tutto
questo tempo. So che mi disapprova.»

«Non te. Forse il tuo stile di vita.»

«Intendi la mia fede. Non c'è niente di male a parlarne. So che Jase si
sente un po'\ldots{} tradito, credo. Come se Simon e io avessimo
ripudiato tutto ciò in cui crede. Ma non è vero. Jason e io non siamo
mai andati nella stessa direzione.»

«Fondamentalmente, sai, è soltanto Jase. Il solito vecchio Jase.»

«E io, sono la solita vecchia Diane?»

A questo non avevo risposta.

Mangiò con evidente appetito e dopo la portata principale ordinammo il
dessert e il caffè. Dissi: «Meno male che sei riuscita a prenderti del
tempo per venire.»

«Meno male che Simon mi ha tolto il guinzaglio?»

«Non intendevo questo.»

«Lo so. Ma in un certo senso è la verità. Simon tende a essere un po'
maniaco del controllo. Gli piace sapere dove sono.»

«Per te è un problema?»

«Intendi se il mio matrimonio è in crisi? No. Non lo è e non lo
permetterei, ma questo non significa che siamo sempre d'accordo su
tutto.» Esitò. «Se ti parlo di queste cose, mi sto confidando con te, va
bene? Non con Jason. Solo con te.»

Annuii.

«Simon è un po' cambiato rispetto a quando lo hai conosciuto tu. Tutti
lo siamo, tutti quelli dei tempi dell'NR. L'NR consisteva nell'essere
giovani e formare una comunità di fede, una sorta di spazio sacro in cui
non dovevamo avere paura gli uni degli altri, dove potersi abbracciare
non solo in senso figurato ma letterale. Credevamo fosse un paradiso
terrestre, ma ci sbagliavamo. Pensavamo che l'AIDS non avesse
importanza, la gelosia non avesse importanza\ldots{} non potevano
averne, perché eravamo arrivati alla fine del mondo. Ma è una
Tribolazione lenta, Ty. La Tribolazione è l'opera di una vita e per lei
dobbiamo essere forti e in salute.»

«Tu e Simon\ldots»

«Oh, stiamo bene.» Sorrise. «E grazie per averlo chiesto, dottor Dupree.
Ma abbiamo perso degli amici per l'AIDS e le droghe. Il movimento era
una corsa sulle montagne russe, amore lungo la salita e dolore lungo la
discesa. Chiunque ne abbia fatto parte ti dirà la stessa cosa.»

Probabilmente era così ma l'unica veterana dell'NR che conoscevo era
Diane stessa. «Gli ultimi anni non sono stati facili per nessuno.»

«Simon ha faticato ad affrontare la situazione. Credeva veramente che
fossimo una generazione benedetta. Una volta mi ha detto che Dio si era
avvicinato tanto all'umanità che sembrava di stare seduti accanto a una
fornace in una notte d'inverno, che praticamente ci si poteva scaldare
le mani nel Regno dei Cieli. Ci sentivamo tutti così, ma per Simon fu
qualcosa di più profondo, quella sensazione riuscì davvero a tirare
fuori il meglio da lui. E quando cominciò ad andare male, quando tanti
nostri amici si ammalarono o finirono nel tunnel delle dipendenze di
vario genere, rimase profondamente colpito. Quello fu anche il periodo
in cui cominciarono a finire i soldi e Simon dovette cercarsi un
lavoro\ldots{} entrambi in realtà. Per qualche anno mi sono barcamenata
tra occupazioni a tempo determinato. Nel frattempo, Simon non è riuscito
a trovare un lavoro nel mondo laico e ora svolge mansioni di portierato
nella nostra chiesa di Tempe, il Jordan Tabernacle, e quando possono lo
pagano\ldots{} sta studiando per il diploma da idraulico.»

«Non è proprio la Terra Promessa.»

«Già, ma sai una cosa? Non credo che debba esserlo. È quello che gli
dico sempre. Forse possiamo sentire l'arrivo del chiliasmo, ma non è
ancora qui\ldots{} dobbiamo ancora giocare gli ultimi minuti della
partita, anche se il risultato è scontato. E forse saremo giudicati
proprio per questo. Dobbiamo giocare come se ne valesse la pena.»

Prendemmo l'ascensore per salire nelle nostre stanze. Diane si fermò un
attimo sulla porta e disse: «In questo momento sto ricordando a me
stessa quanto sia bello parlare con te. Eravamo piuttosto loquaci,
ricordi?»

Ci confidavamo le nostre paure attraverso un casto mezzo telefonico.
Intimità a distanza. Aveva sempre preferito fare così. Annuii.

«Forse possiamo farlo ancora» disse. «Forse qualche volta potrei
chiamarti dall'Arizona.»

Ovviamente sarebbe stata lei a chiamare me perché Simon non avrebbe
gradito che chiamassi io. Era sottinteso, come la natura del rapporto
che mi stava proponendo. Sarei stato il suo amico platonico. Una persona
innocua con cui confidarsi nei momenti difficili, come l'amico gay della
protagonista di un dramma sentimentale in un cinema multisala. Avremmo
chiacchierato. Ci saremmo confidati. Nessuno avrebbe sofferto.

Non era proprio ciò che avrei voluto o di cui sentivo il bisogno, ma non
potevo dirlo di fronte allo sguardo entusiasta e leggermente perso che
mi stava lanciando. Dissi: «Sì, certo.»

E così sorrise, mi abbracciò e mi lasciò in corridoio.

Mi alzai più tardi di quanto avrei dovuto, mi presi cura della mia
dignità ferita, nascosto fra i rumori e le risate delle camere
adiacenti, pensando a tutti gli scienziati e gli ingegneri della
Perielio, del JPL e del Kennedy, a tutta quella gente dei quotidiani e
ai giornalisti della TV che guardavano le luci dei riflettori guizzare
sui razzi lontani. Tutti noi facevamo la nostra parte lì, sullo
striscione del traguardo della storia umana. Facevamo quello che ci si
aspettava da noi, ce la giocavamo come se ne valesse la pena.

ooo

Jason arrivò a mezzogiorno del giorno dopo, dieci ore prima dell'inizio
programmato della prima ondata di lanci. Le condizioni meteo erano
radiose e calme, un buon auspicio. Tra tutti i siti di lancio al mondo,
l'unico che ricevette un chiaro invito a non operare fu il complesso
Kourou dell'Agenzia Spaziale Europea nella Nuova Guinea francese, chiuso
a causa di un violento temporale di marzo. (I microorganismi
dell'Agenzia Spaziale Europea avrebbero subito un ritardo di uno o due
giorni\ldots{} o mezzo milione di anni, secondo il tempo dello Spin.)

Jase venne direttamente in camera mia, dove io e Diane lo stavamo
aspettando. Indossava una giacca a vento di plastica da due soldi e un
cappellino dei Marlins tirato giù fino a coprirgli gli occhi per
camuffarsi tra i cronisti che soggiornavano nell'albergo. «Tyler,»
esclamò quando aprii la porta, «mi dispiace. Se avessi potuto sarei
venuto.»

Il funerale. «Lo so.»

«Belinda Dupree era la parte migliore della Grande Casa. Dico sul
serio.»

«Ti ringrazio» dissi, e mi spostai per farlo passare.

Diane attraversò la stanza con un'espressione guardinga. Jason chiuse la
porta dietro di sé, senza sorridere. Erano a un metro di distanza l'uno
dall'altra e si adocchiavano. Il silenzio era pesante. Jason lo ruppe.

«Quel colletto,» disse, «ti fa assomigliare a un banchiere vittoriano. E
poi dovresti mettere su qualche chilo. È così difficile scroccare un
pasto laggiù nel paese delle mucche?»

Diane rispose: «Più cactus che mucche, Jase.»

E scoppiarono a ridere e si gettarono l'uno nelle braccia dell'altra.

ooo

Ci appostammo fuori sul terrazzo appena fece buio, portammo fuori delle
sedie comode e ordinammo un vassoio di crudités (fu Diane a scegliere)
dal servizio in camera. La notte era scura, come ogni altra notte senza
stelle avvolta dallo Spin, ma le piattaforme di lancio erano illuminate
da giganteschi fari i cui riflessi ballavano sulle onde che rollavano
dolcemente.

Da alcune settimane ormai, Jason era in cura da un neurologo. La
diagnosi dello specialista era stata la stessa che avevo già fatto io:
Jason era affetto da una grave forma di sclerosi multipla che non
rispondeva alle cure e il cui unico trattamento efficace consisteva in
una terapia di farmaci palliativi. Di fatto il neurologo aveva voluto
sottoporre il caso di Jason al Centro per il Controllo Malattie
inserendolo come parte dei loro studi in via di sviluppo su ciò che
alcuni chiamavano SMA: sclerosi multipla atipica. Jase lo aveva
minacciato o corrotto per fare in modo che si togliesse dalla testa
quell'idea. E per il momento, almeno, il nuovo cocktail lo stava tenendo
in fase di remissione. Era operativo e riusciva a muoversi come al
solito. Qualunque sospetto avesse covato Diane venne ben presto fugato.

Per celebrare i lanci aveva portato con sé una bottiglia di autentico
champagne francese molto costoso. «Avremmo potuto ottenere dei posti in
zona VIP» dissi a Diane. «Le tribune fuori dell'Edificio Assemblaggio
Veicoli. Gomito a gomito con il Presidente Garland.»

«La vista da qui è ugualmente buona» disse Jason. «Meglio. Qui non siamo
oggetti di scena in una foto di gruppo.»

«Non ho mai incontrato un Presidente» disse Diane.

Il cielo, ovviamente, era buio, ma la TV nella camera d'albergo (avevamo
alzato il volume per sentire il conto alla rovescia) parlava della
barriera dello Spin e Diane guardava verso il cielo, come se quel
coperchio che aveva racchiuso il mondo fosse diventato miracolosamente
visibile. Jason notò che inclinava la testa. «Non dovrebbero chiamarla
barriera» disse. «Nessuna delle riviste specializzate la chiama più
così.»

«Eh? E come la chiamano?»

Si schiarì la voce. «Una ``strana membrana''.»

«Oh no.» Diane rise. «No, è orribile. Non si può sentire. Sembra una
malattia ginecologica.»

«Sì, ma ``barriera'' non è corretto. È più uno strato limite. Non è una
linea che si attraversa. Acquisisce gli oggetti in maniera selettiva e
li accelera verso l'universo esteriore. Il processo è più simile
all'osmosi che, diciamo, allo sfondamento di uno steccato. Ergo,
membrana.»

«Avevo dimenticato come fosse parlare con te, Jase. Leggermente
surreale.»

«Zitti» dissi a entrambi. «Ascoltate.»

Adesso la TV si era collegata in diretta con la NASA, una voce insipida
dal Controllo della Missione faceva il conto alla rovescia. Trenta
secondi. C'erano dodici razzi sulle piattaforme, riforniti e pronti al
lancio. Dodici lanci simultanei, un atto che un'agenzia spaziale meno
ambiziosa avrebbe un tempo ritenuto impraticabile e totalmente
pericoloso. Ma vivevamo in tempi più audaci o più disperati.

«Perché devono partire tutti insieme?» chiese Diane.

«Perché\ldots» cominciò Jase; poi disse: «No. Aspetta. Guarda.»

Venti secondi. Dieci. Jase si alzò in piedi e si sporse dal corrimano
del terrazzo. I terrazzi dell'albergo erano assediati. La spiaggia era
assediata. Migliaia di teste e obiettivi puntati nella stessa direzione.
Le stime successive stabilirono che la folla dentro e fuori del Cape si
aggirava sui due milioni di persone. Secondo i bollettini della polizia
quella notte vennero rubati più di cento portafogli. Ci furono due
accoltellamenti fatali, quindici tentate aggressioni e un parto
prematuro. (La neonata, una bambina di un chilo e ottocento grammi,
venne data alla luce su un tavolo poggiato su cavalletti nella Casa
Internazionale delle Frittelle al Cocoa Beach.)

Cinque secondi. La TV della camera d'albergo si zittì. Per un attimo non
ci fu alcun rumore tranne il ronzio e il sibilo delle attrezzature
fotografiche.

Poi l'oceano venne incendiato dalla luce delle fiamme fino
all'orizzonte.

Nessuno di quei razzi da solo avrebbe impressionato la folla, neppure al
buio, ma questa non era una colonna di fuoco, erano cinque, sette,
dieci, dodici. Le rampe marittime si stagliarono brevemente controluce
come grattacieli scheletrici, per perdersi subito dopo tra le nuvole
delle acque oceaniche vaporizzate. Dodici pilastri di fuoco bianco,
separati da vari chilometri, ma compressi dalla prospettiva,
artigliarono il cielo diventato color indaco per la combinazione delle
luci. La folla sulla spiaggia cominciò a festeggiare e il suono si fuse
con il rumore dei propulsori a combustibile solido che martellavano per
l'altitudine, una vibrazione che stringeva il cuore come l'estasi o il
terrore. Ma non era solamente lo spettacolo in sé che stavamo
festeggiando. Quasi sicuramente ognuno tra quei due milioni di persone
aveva già visto il lancio di un razzo, almeno in televisione, e per
quanto quella multipla ascesa fosse imponente e rumorosa, era
straordinaria soprattutto per il suo intento, per l'idea che l'aveva
determinata. Non stavamo solo piantando la bandierina della vita
terrestre su Marte, stavamo sfidando lo Spin stesso.

I razzi salivano. (E sullo schermo rettangolare della TV, quando lanciai
un'occhiata attraverso la porta del terrazzo, razzi simili s'inserirono
nella nuvolosa luce diurna a Jiuquan, Svobodnyj, Bayqoñyr, Xichang.)
L'impetuosa luce orizzontale divenne obliqua e cominciò ad affievolirsi
mentre la notte tornava di fretta dal mare. Il rumore si esaurì nella
sabbia, nel cemento e nell'acqua salata surriscaldata. Immaginai di
riuscire a sentire l'odore acre dei fuochi d'artificio che giungeva a
riva con la marea, simile al tanfo tremendo ma gradevole delle candele
romane.

Un migliaio di macchine fotografiche ronzarono come grilli in fin di
vita e poi tacquero.

I festeggiamenti durarono, in un modo o nell'altro, fino all'alba.

ooo

Rientrammo dentro, tirammo le tende per non vedere l'oscurità
sconfortante e stappammo lo champagne. Guardammo i notiziari
dall'estero. A parte il ritardo francese dovuto alla pioggia, ogni
lancio era avvenuto con successo. Un'armata batterica era in rotta su
Marte.

«Quindi, perché devono \emph{proprio} partire tutti contemporaneamente?»
chiese di nuovo Diane.

Jason le lanciò un lungo sguardo riflessivo. «Perché vogliamo che
arrivino a destinazione pressappoco nello stesso momento. Cosa che non è
così facile come sembra. Devono penetrare nella membrana dello Spin
quasi simultaneamente, altrimenti potrebbero uscirne a distanza di anni
o secoli. Non è una criticità per questi carichi anaerobici ma ci
esercitiamo per quando sarà davvero importante.»

«Anni o \emph{secoli}? Com'è possibile?»

«La natura dello Spin, Diane.»

«Va bene, ma secoli?»

Girò la sedia mettendosi di fronte a lei, accigliato. «Voglio provare a
comprendere l'entità della tua ignoranza\ldots»

«È solo una domanda, Jase.»

«Conta un secondo.»

«Come?»

«Guarda l'orologio e conta un secondo. No, lo faccio io. Un.» Pausa.
«Secondo. È chiaro?»

«Jason\ldots»

«Abbi pazienza. Capisci il coefficiente dello Spin?»

«Più o meno.»

«Più o meno non basta. Un secondo terrestre corrisponde a 3,17 anni di
tempo dello Spin. Tienilo a mente. Se uno dei nostri razzi entra nella
membrana dello Spin un solo secondo dopo gli altri, raggiunge l'orbita
con più di tre anni di ritardo.»

«Solo perché non riesco a citare i numeri\ldots»

«Sono numeri importanti, Diane. Supponiamo che la nostra flottiglia sia
appena emersa dalla membrana, proprio adesso, \emph{ora}\ldots» fece un
segno nell'aria con il dito. «Un \emph{secondo}, è qui e non c'è più.
Per la flottiglia questi erano tre anni e spiccioli. Un secondo fa erano
nell'orbita terrestre. Ora hanno consegnato il carico sulla superficie
di Marte. Intendo \emph{ora}, Diane, letteralmente \emph{ora}. È già
successo, è fatta. Fai trascorrere un minuto sul tuo orologio. Questi
sono approssimativamente centonovant'anni su un orologio esterno.»

«Sono tanti, certo, ma non puoi modificare un pianeta in duecento anni,
o no?»

«Dunque, ora l'esperimento dura da duecento anni di Spin. Proprio ora,
mentre parliamo, tutte le colonie batteriche che sono sopravvissute al
viaggio si saranno riprodotte su Marte da due secoli. Fra un'ora,
saranno lì da undicimilaquattrocento anni. Domani a quest'ora si saranno
moltiplicati da circa duecentosettantaquattromila anni.»

«Va bene, Jase. Mi sono fatta un'idea.»

«La prossima settimana a quest'ora, 1,9 milioni di anni.»

«Va bene.»

«Un mese, 8,3 milioni di anni.»

«Jason\ldots»

«L'anno prossimo a quest'ora, cento milioni di anni.»

«Sì, ma\ldots»

«Sulla Terra, cento milioni di anni è circa il tempo trascorso
dall'inizio del Cretacico superiore al tuo ultimo compleanno. Cento
milioni di anni sono una quantità di tempo sufficiente perché quei
microorganismi possano pompare il biossido di carbonio fuori dai
depositi sulla crosta, filtrare l'azoto dai nitrati, depurare gli ossidi
dalla regolite e arricchirli morendo in grosse quantità. Tutta
l'anidride carbonica liberata è un gas a effetto serra. L'atmosfera
diventa più densa e più calda. Tra un anno invieremo un'altra armata di
organismi che respirano e loro cominceranno il ciclo dall'anidride
carbonica all'ossigeno libero. Un altro anno - oppure appena la firma
spettroscopica proveniente dal pianeta sarà quella giusta - e
introdurremo l'erba, le piante e altri organismi complessi. E quando
tutto ciò si stabilizzerà in una sorta di ecologia planetaria
grossolanamente omeostatica, manderemo gli esseri umani. Sai cosa vuol
dire?»

«Dimmelo tu» disse Diane sgarbatamente.

«Vuol dire che nel giro di cinque anni ci sarà una fiorente
civilizzazione umana su Marte. Fattorie, fabbriche, strade, città\ldots»

«Esiste una parola greca per indicare tutto questo, Jase.»

«Ecopoiesi.»

«Io stavo pensando a ``hybris''.»

Sorrise. «Mi preoccupo per un sacco di cose. Ma offendere gli dei non è
una di queste.»

«Oppure offendere gli Ipotetici?»

Questo lo bloccò. Si appoggiò alla sedia e prese un sorso di champagne,
ormai un po' sgassato, dal suo bicchiere dell'albergo.

«Non ho paura di offenderli» disse alla fine. «Al contrario. Temo che
forse stiamo facendo esattamente ciò che vogliono che facciamo.»

Ma non dette ulteriori spiegazioni e Diane era impaziente di cambiare
argomento.

ooo

Accompagnai Diane a Orlando il giorno dopo, dove la attendeva un volo
per Phoenix.

Negli ultimi giorni era diventato fin troppo ovvio che non avremmo
discusso, menzionato o accennato in alcun modo all'intimità fisica che
c'era stata quella notte nel Berkshire prima che si sposasse con Simon.
L'unica ammissione fu evitare l'argomento con deviazioni farraginose.
Quando ci abbracciammo (in modo casto) nell'area davanti all'ingresso
della sicurezza aeroportuale mi disse ``ti chiamo'' e sapevo che lo
avrebbe fatto - Diane non faceva molte promesse ma era scrupolosa nel
mantenerle - ma ero altrettanto consapevole del tempo che era trascorso
dall'ultima volta che l'avevo vista e del tempo che sarebbe
inevitabilmente passato prima che potessi rivederla: non il tempo dello
Spin, ma qualcosa di altrettanto corrosivo e famelico. Aveva delle
pieghe sugli angoli degli occhi e della bocca, non molto diverse da
quelle che vedevo allo specchio ogni mattina.

Incredibile, pensai, con che velocità eravamo diventati come quelle
persone che non si conoscono molto bene tra di loro.

ooo

Durante la primavera e l'estate di quell'anno ci furono altri lanci,
pacchetti di sorveglianza che passavano mesi nell'orbita alta della
Terra e tornavano portando immagini visive e spettrografiche di
Marte\ldots{} istantanee dell'ecopoiesi.

I primi risultati furono ambigui: un modesto incremento dell'anidride
carbonica atmosferica che poteva essere un effetto collaterale del
riscaldamento solare. Marte continuava a essere un mondo freddo e
inospitale sotto ogni aspetto plausibile. Jason ammise che persino i
GEMO - gli organismi geneticamente modificati di Marte che costituivano
la gran parte dell'inseminazione iniziale - forse non si erano ben
adattati ai livelli degli ultravioletti della luce diurna non filtrata
del pianeta e alla regolite carica di agenti ossidanti.

Ma per la metà dell'estate notammo prove spettrografiche evidenti di
un'attività biologica. C'era una quantità maggiore di vapore acqueo in
un'atmosfera più densa, maggiore quantità di metano, etano e ozono, e
persino un leggerissimo ma rilevabile aumento di azoto libero.

Per Natale quei cambiamenti, per quanto ancora impercettibili, avevano
drasticamente superato ciò che si poteva attribuire al riscaldamento
solare tanto che non vi erano più dubbi, Marte era diventato un pianeta
vivente.

Le piattaforme di lancio vennero allestite nuovamente, nuovi carichi di
vita microbica coltivata e impacchettata. Quell'anno negli Stati Uniti
il due per cento del prodotto interno lordo fu destinato al campo
aerospaziale relativo allo Spin - essenzialmente al programma Marte - e
più o meno la stessa percentuale anche negli altri paesi
industrializzati.

ooo

A febbraio Jason ebbe una ricaduta. Si svegliò non riuscendo a mettere a
fuoco. Il suo neurologo calibrò la terapia farmacologica e gli
prescrisse una benda sugli occhi come soluzione temporanea. Jase si
riprese in fretta ma rimase lontano dal lavoro per quasi tutta la
settimana.

Diane fu di parola. Cominciò a chiamarmi almeno una volta al mese, di
solito di più, spesso a tarda sera quando Simon dormiva dall'altra parte
del loro piccolo appartamento. Vivevano in poche stanze sopra un negozio
di libri di seconda mano a Tempe, il massimo che potevano fare in base
ai guadagni di Diane e al reddito saltuario che Simon portava a casa dal
Jordan Tabernacle. Quando faceva caldo, riuscivo a sentire in sottofondo
il ronzio di un condizionatore evaporativo; d'inverno una radio a volume
basso per camuffare il suono della sua voce.

La invitai a tornare in Florida per la serie di lanci successiva ma
ovviamente non poteva: aveva impegni di lavoro, quel fine settimana
avevano ospiti della chiesa a cena, Simon non avrebbe capito. «Simon sta
attraversando una leggera crisi spirituale. Sta cercando di occuparsi
della questione del Messia\ldots»

«Esiste una questione del Messia?»

«Dovresti leggere i giornali» disse Diane. Probabilmente sovrastimava la
presenza di certi dibattiti religiosi sulla stampa tradizionale, almeno
in Florida; forse a Ovest era diverso. «La vecchia corrente NR credeva
in una Parusia senza Cristo. Era questo che ci distingueva.» Quello e la
loro inclinazione alla nudità pubblica, pensai. «I primi scrittori,
Ratel e Greengage, vedevano nello Spin la realizzazione diretta delle
profezie delle Scritture - che significava che la profezia stessa veniva
ridefinita, riconfigurata dagli eventi storici. Non doveva esserci una
vera e propria Tribolazione e nemmeno una seconda venuta di Cristo in
carne e ossa. Tutta quella roba nei Tessalonicesi, nei Corinzi e nella
Rivelazione poteva essere reinterpretata o ignorata, perché lo Spin era
un intervento genuino di Dio nella storia dell'uomo - un miracolo
tangibile che sostituiva le Scritture. Era ciò che ci aveva liberato
perché costruissimo un Regno sulla Terra. All'improvviso eravamo
responsabili del nostro stesso chiliasmo.»

«Non credo di seguirti.» In realtà mi ero già perso alla parola
``Parusia''.

«Significa\ldots{} be', ciò che conta davvero è che il Jordan
Tabernacle, la nostra piccola chiesa, ha rinunciato ufficialmente a
tutta la dottrina dell'NR, anche se metà della congregazione è formata
da vecchi adepti come me e Simon. Così all'improvviso sono iniziate
tutte queste discussioni sulla Tribolazione e su quanto lo Spin collimi
con le profezie bibliche. La gente si è arroccata sulle idee più
svariate. Esiste un Anticristo e, se sì, dove si trova? L'Assunzione
avviene prima, durante o dopo la Tribolazione? Questioni del genere.
Forse ti sembreranno battibecchi da niente ma la posta spirituale in
gioco è alta e quelli che discutono su questi argomenti sono persone a
cui teniamo, amici nostri.»

«Tu da che parte stai?»

«Personalmente?» Rimase in silenzio e poi eccola di nuovo, la musica
della radio che mormorava alle sue spalle, un annunciatore dalla voce
soporifera che distribuiva le notizie della notte agli insonni. Le
ultime sulla sparatoria a Mesa. Parusia o non Parusia. «Posso dire di
essere combattuta. Non so a cosa credo. A volte sento la mancanza dei
vecchi tempi. Costruire il paradiso giorno per giorno. È un po'
come\ldots»

Si bloccò. Adesso c'era un'altra voce a moltiplicare per due il mormorio
e l'interferenza della radio: «Diane? Sei ancora in piedi?»

«Scusa» sussurrò. Simon in perlustrazione. Era ora di interrompere la
nostra tresca telefonica, il suo atto d'infedeltà senza contatto. «Ci
sentiamo presto.»

Se n'era andata prima ancora che potessi salutarla.

ooo

La seconda serie di lanci dei semi si svolse in maniera impeccabile come
la prima. I media affollarono di nuovo Canaveral, ma stavolta guardai
l'evento su un megaschermo presso l'auditorio della Perielio, un lancio
in pieno giorno che sparpagliò gli aironi nel cielo sopra Merritt Island
come fossero coriandoli lucenti.

Seguì un'altra estate di attesa. L'Agenzia Spaziale Europea lanciò in
orbita una serie di telescopi di ultima generazione e degli
interferometri. I dati che raccolsero furono addirittura migliori e più
chiari rispetto all'anno precedente. Verso settembre tutti gli uffici
della Perielio erano tappezzati d'immagini ad alta risoluzione del
nostro successo. Ne incorniciai uno per la sala d'attesa
dell'infermeria. Era un render a colori compositi di Marte, che mostrava
l'Olympus Mons delineato dal gelo o dal ghiaccio e segnato da recenti
canali di drenaggio, con la nebbia che scorreva come acqua attraverso le
Valles Marineris e capillari verdi che serpeggiavano sul Solis Lacus.
Gli altopiani meridionali della Terra Sirenum erano deserti immobili ma
i crateri da impatto della regione erano stati erosi dal clima divenuto
più umido e ventoso fin quasi a diventare invisibili.

Il contenuto d'ossigeno dell'atmosfera crebbe e diminuì per alcuni mesi
come l'oscillazione della popolazione degli organismi aerobici, ma verso
dicembre superò i venti millibar e si stabilizzò. Marte stava scoprendo
il proprio equilibrio attraverso un miscuglio potenzialmente caotico di
gas serra in aumento, un ciclo idrologico instabile e nuovi cicli
biogeochimici di retroazione.

La catena di successi fece bene a Jason. Continuava a essere in fase di
remissione ed era felicemente indaffarato, quasi in modo terapeutico. Se
c'era qualcosa che lo aveva sgomentato, era il suo ruolo di genio
iconico della Fondazione Perielio, o perlomeno la sua celebrità
scientifica, immagine simbolo della trasformazione di Marte. Tutto ciò
era dovuto più agli sforzi di E.D. che di Jason: E.D. sapeva che
l'opinione pubblica voleva che la Perielio avesse un volto umano,
preferibilmente giovane e intelligente, ma non minaccioso. Per questo
aveva spinto Jase di fronte alle telecamere sin dai giorni in cui la
Perielio era una lobby aerospaziale. Jase si era rassegnato a farlo -
era un divulgatore bravo, paziente e ragionevolmente fotogenico - ma
odiava quella prassi e avrebbe preferito stare in disparte piuttosto che
vedersi in televisione.

Quello fu l'anno dei primi voli PNE senza equipaggio, che Jase osservava
con particolare attenzione. Questi erano i veicoli che avrebbero
trasportato gli esseri umani su Marte e, diversamente dai portatori di
semi, al confronto molto semplici, i veicoli PNE erano la nuova
tecnologia. PNE stava a indicare ``propulsione nucleare-elettrica'':
reattori nucleari in miniatura che alimentavano motori ionici molto più
potenti di quelli che guidavano le navicelle con i semi e abbastanza
potenti da permettere carichi enormi. Ma per portare in orbita questi
giganti servivano i razzi più grandi che la NASA avesse mai lanciato.
Secondo Jason serviva un'operazione di ``ingegneria eroica'', ovvero
eroicamente cara. L'enorme costo della missione aveva cominciato a far
scattare segnali di allerta persino all'interno di un Congresso
ampiamente solidale, ma la sequenza di notevoli successi teneva a freno
il dissenso. Jason temeva che anche il minimo fallimento avrebbe
spostato i termini dell'equazione.

Poco dopo Capodanno un veicolo PNE di prova non riuscì a riportare il
suo pacchetto di rientro con i dati di analisi e si ipotizzò che fosse
rimasto in panne in orbita. A Capitol Hill, una cricca di stati
fiscalmente ultraconservatori che non avevano investito in maniera
significativa nel comparto aerospaziale condusse una campagna di
biasimo, ma gli amici di E.D. al Congresso misero a tacere le obiezioni
e un test positivo effettuato una settimana dopo mise fine a ogni
controversia. Comunque, disse Jason, avevamo schivato un proiettile.

Diane aveva seguito il dibattito ma lo considerava insignificante.
«Quello di cui si \emph{deve} preoccupare Jase,» disse, «è ciò che
questa faccenda di Marte sta facendo al mondo. Finora la stampa è stata
tutta favorevole, giusto? Sono tutti entusiasti, tutti vogliamo qualcosa
che ci rassicuri a proposito della - non so bene come definirla -
\emph{potenza} della specie umana. Ma l'euforia alla fine svanirà e nel
frattempo la gente sta diventando molto ferrata sulla natura dello
Spin.»

«Cosa c'è di male?»

«Il progetto Marte potrebbe fallire o non essere all'altezza delle
aspettative, ecco cosa. E non solo perché la gente sarà amareggiata.
Hanno osservato la trasformazione di un intero pianeta e ora hanno un
righello per misurare lo Spin. Ovviamente mi riferisco al suo potere
insensato. Ma lo Spin non è un semplice fenomeno astratto. Avete fatto
in modo che la gente guardasse la bestia negli occhi e questo per voi è
un buon risultato, suppongo. Ma se il vostro progetto non funziona,
tutto quel coraggio verrà strappato via di nuovo e a quel punto sarà
peggio perché ormai avranno \emph{visto} la bestia. E non vi
compatiranno per aver fallito, Tyler, perché saranno ancora più
spaventati di prima.»

Citai a memoria la poesia di Housman che mi aveva insegnato tanto tempo
prima: ``\emph{All'infante non è occorso / che se l'è mangiato
	l'orso}.''

«\emph{L'infante} inizia a capire» disse. «Forse è un modo di intendere
la Tribolazione.»

Forse era così. Certe notti, quando non riuscivo a prendere sonno,
pensavo agli Ipotetici, chiunque o qualunque cosa fossero. C'era un solo
fatto che li riguardava e che era davvero importante e chiaro: non solo
erano stati in grado di racchiudere la Terra in questa\ldots{} strana
membrana, ma erano attorno a noi - a possederci, a regolare il nostro
pianeta e lo scorrere del tempo - da quasi due miliardi di anni.

Niente che fosse anche solo lontanamente umano poteva essere così
paziente.

ooo

Il neurologo di Jason mi riferì di uno studio pubblicato quell'inverno
sul JAMA, il giornale dell'Associazione Medica Americana. I ricercatori
della Cornell avevano scoperto un marcatore genetico per i casi acuti di
sclerosi multipla farmaco-resistente. Il neurologo - un tipo grosso e
geniale della Florida di nome David Malmstein - aveva fatto delle
ricerche sul DNA di Jason e vi aveva trovato una sequenza sospetta. Gli
chiesi cosa significasse.

«Significa che possiamo personalizzare la sua terapia in modo un po' più
specifico. Significa anche che non possiamo in alcun modo assicurare una
remissione permanente come quella di un tipico paziente con sclerosi
multipla.»

«Pare che la sua fase di remissione stia durando da quasi un anno. Non
ti sembra troppo?»

«I suoi sintomi sono sotto controllo, tutto qui. La sclerosi multipla
atipica continua a bruciare, come fosse un incendio in un giacimento di
carbone. Arriverà il momento in cui non riusciremo a compensarla.»

«Il punto di non ritorno.»

«Diciamo così.»

«Per quanto tempo ancora potrà sembrare normale?»

Malmstein si bloccò. «Sai, è la stessa cosa che mi ha chiesto Jason.»

«Cosa gli hai risposto?»

«Che non sono un veggente. Che la sclerosi multipla atipica è una
malattia senza un'eziologia ben definita. Che il corpo umano ha il
proprio calendario.»

«Immagino che non abbia gradito la risposta.»

«No. Ha fatto sentire il suo dissenso. Ma è vero, potrebbe andarsene in
giro senza alcun sintomo per i prossimi dieci anni. Oppure potrebbe
ritrovarsi su una sedia a rotelle entro la fine della settimana.»

«Gli hai detto \emph{questo}?»

«Una versione più delicata e gentile. Non voglio che perda le speranze.
Ha uno spirito combattivo e questo conta moltissimo. Onestamente, sono
dell'opinione che starà bene nel breve periodo\ldots{} due anni, cinque
o forse di più. Poi, non ci sono più garanzie. Mi spiace, avrei voluto
che la prognosi fosse migliore.»

Non raccontai a Jase che avevo parlato con Malmstein, ma vidi come nelle
settimane successive intensificò i suoi ritmi di lavoro calcolando i
propri successi contro il tempo e la mortalità, non quella del mondo ma
la sua.

ooo

Il ritmo dei lanci, per non parlare dei costi, cominciò a
intensificarsi. L'ultima ondata di lanci di semi (l'unica a trasportare,
in parte, semi veri e propri) fu a marzo, due anni dopo che io, Jase e
Diane avevamo visto una decina di razzi simili partire dalla Florida
verso quello che al tempo era un pianeta sterile.

Lo Spin ci aveva dato lo stimolo necessario a una lunga ecopoiesi. Ora
che avevamo lanciato i semi delle piante complesse, però, la tempistica
diventava cruciale. Se avessimo aspettato troppo Marte avrebbe potuto
evolversi fuori dal nostro controllo: una varietà di grano commestibile
dopo un milione di anni di evoluzione allo stato brado avrebbe potuto
non somigliare per nulla alla sua forma ancestrale, sarebbe potuta
crescere immangiabile o addirittura velenosa.

Questo voleva dire che i satelliti sonda dovevano essere lanciati solo
alcune settimane dopo l'armata di semi. Se i risultati fossero parsi
promettenti, le navicelle PNE con l'equipaggio sarebbero state inviate
subito dopo.

Ricevetti un'altra chiamata notturna da Diane la sera dopo il decollo
dei satelliti sonda (i cui pacchetti dati erano stati recuperati nel
giro di poche ore ed erano in viaggio verso il JPL di Pasadena per
essere analizzati). Mi sembrava stressata e dovetti insistere un bel po'
per farmi raccontare che era stata congedata dal lavoro almeno fino a
giugno. Lei e Simon erano finiti nei guai per l'affitto arretrato. Non
poteva chiedere i soldi a E.D. ed era impossibile parlare con Carol.
Stava cercando di trovare il coraggio di parlarne con Jase ma l'idea di
umiliarsi non la entusiasmava affatto.

«Di quanti soldi stiamo parlando, Diane?»

«Tyler, non intendevo\ldots»

«Lo so. Non mi hai chiesto nulla. Te lo sto proponendo io.»

«Be'\ldots{} questo mese, anche cinquecento dollari farebbero una grossa
differenza.»

«Suppongo che il patrimonio dell'idraulico si sia prosciugato.»

«Il fondo fiduciario di Simon è finito. Ci sono ancora i soldi di
famiglia, ma la sua famiglia non gli parla più.»

«Non è che se ne accorgerà, se ti mando un assegno?»

«Non gli farebbe piacere. Pensavo di raccontargli di aver ritrovato una
vecchia polizza assicurativa e di averla incassata. Una cosa così. Il
genere di bugia che non conta come peccato. Spero.»

«Voi due siete ancora al solito indirizzo di Collier Street?» Dove ogni
anno spedivo una cartolina di Natale educata e neutra, per riceverne una
in risposta: generiche scene di neve firmate: ``Simon e Diane Townsend,
Dio ti benedica!''

«Sì,» disse, e poi, «grazie, Tyler. Grazie mille. Non sai quanto sia
mortificante.»

«Sono tempi duri per un sacco di gente.»

«Invece, tu stai bene?»

«Sì, tutto a posto.»

Le spedii sei assegni ognuno dei quali postdatato al quindici del mese,
valevano metà del mio stipendio annuo, senza sapere se questo avrebbe
rinsaldato la nostra amicizia o l'avrebbe avvelenata. O se ne valesse la
pena.

ooo

I dati d'indagine evidenziarono un mondo comunque più secco della Terra
ma segnato da laghi, che sembravano raffinati intarsi turchesi su un
disco di rame; un pianeta con leggere fasce di nuvole vorticanti e
temporali che lasciavano cadere la pioggia sui pendii esposti al vento
di antichi vulcani, e alimentavano i bacini fluviali e i limacciosi
delta, pianeggianti e verdi come prati di periferia.

Gli enormi razzi vennero riforniti sulle loro piattaforme. Nelle
infrastrutture di lancio e nei cosmodromi in giro per il mondo quasi
ottocento esseri umani si arrampicarono sulle rampe per rinchiudersi in
cavità grandi quanto una credenza e affrontare un destino incerto. Le
arche PNE, contenute nella cuspide di quei razzi, alloggiavano (oltre
agli astronauti) i grembi d'acciaio da cui, con un po' di buona sorte,
sarebbero stati travasati embrioni di pecore, bovini, maiali e capre; i
semi di diecimila piante; larve di api e altri preziosi insetti; decine
di carichi biologici simili che avrebbero potuto sopravvivere, oppure
no, al viaggio e alle difficoltà della rigenerazione; archivi compatti
dell'essenza del sapere umano, sia in digitale (accompagnato dagli
strumenti per la lettura) che in cartaceo (stampato in modo fitto); e
componenti e forniture per semplici rifugi, generatori a energia solare,
serre, depuratori per l'acqua e ospedali da campo rudimentali. Nella
migliore delle ipotesi tutte queste navicelle sarebbero arrivate
all'incirca sulle stesse pianure equatoriali nell'arco di vari anni a
seconda del loro transito attraverso la membrana dello Spin. Alla
peggio, se anche una sola nave fosse arrivata ragionevolmente intatta,
sarebbe stata in grado di sostentare il proprio equipaggio durante il
periodo di acclimatazione.

Ero ancora una volta nell'auditorium della Perielio insieme a tutti
quelli che non avevano risalito la costa per andare a vedere l'evento di
persona. Sedevo in prima fila accanto a Jason. Allungavamo il collo per
vedere le trasmissioni video provenienti dalla NASA, una spettacolare
ripresa a campo lungo delle piattaforme di lancio in alto mare, isole
d'acciaio collegate da immensi ponti ferroviari, dieci giganteschi razzi
Prometheus (chiamati ``Prometheus'' dalla Boeing e dalla
Lockheed-Martin; i russi, i cinesi e gli europei usavano lo stesso
modello ma chiamavano e dipingevano i loro razzi in modi differenti)
inondati di luci e posizionati come i paletti di una staccionata
verniciata di bianco che si estendeva nel blu dell'Atlantico. Molto era
stato sacrificato per quel momento: tasse e patrimoni, litorali e
barriere coralline, carriere e vite. (Ai piedi di ciascuna torre fuori
Canaveral c'era una targa con sopra incisi i nomi dei quindici operai
edili che erano morti durante il montaggio delle costruzioni.) Jason
batteva il piede a ritmo concitato mentre il conto alla rovescia si
svuotava fino all'ultimo minuto. Lo osservavo chiedendomi se tutto ciò
fosse sintomatico, ma Jason se ne accorse e si avvicinò per sussurrarmi
all'orecchio: «Sono solamente un po' nervoso. Tu no?»

C'erano già stati dei problemi. In tutto il mondo, ottanta di quegli
enormi razzi erano stati assemblati e preparati per il lancio
sincronizzato di quella sera. Ma erano progettazioni nuove, non messe a
punto completamente. Quattro lanci erano stati annullati prima del via
per problemi tecnici. Su tre stavano tenendo in sospeso il conteggio -
per un'operazione che doveva essere sincronizzata in tutto il mondo -
per i soliti motivi: tubi di rifornimento pericolosi, intoppi col
software. Tutto ciò era inevitabile e durante la pianificazione era
stato preso in considerazione, ma parve comunque malaugurante.

Tante cose da fare in pochissimo tempo. Stavolta non stavamo
trapiantando biologia ma la storia dell'uomo, e la storia dell'uomo,
aveva detto Jase, bruciava come il fuoco rispetto alla ruggine lenta
dell'evoluzione. (Quando eravamo molto più giovani, dopo lo Spin ma
prima che lasciasse la Grande Casa, Jase faceva un trucchetto da quattro
soldi per dimostrare questa idea. «Allarga le braccia,» diceva,
«stendile bene all'altezza delle spalle» e quando ti aveva messo nella
giusta posizione a croce proseguiva, «immagina che la storia della Terra
sia lunga dall'indice della tua mano sinistra fino a quello della
destra, passando per il cuore. Lo sai quale sarebbe la storia
dell'\emph{uomo}? La storia dell'uomo sarebbe l'unghia dell'indice della
destra. E nemmeno tutta l'unghia. Solo quella piccola parte bianca. La
parte che tagli quando è troppo lunga. Quella è la scoperta del fuoco e
l'invenzione della scrittura e Galileo e Newton e lo sbarco sulla Luna e
l'undici settembre e la settimana scorsa e stamattina. Rispetto
all'evoluzione siamo dei neonati. Rispetto alla geologia esistiamo a
malapena».)

Poi la voce della NASA annunciò: «Accensione» e Jason risucchiò l'aria
tra i denti e girò la testa a metà mentre nove razzi su dieci, tubi cavi
pieni di liquido esplosivo e più alti dell'Empire State Building,
esplodevano verso il cielo contro ogni logica di gravità e inerzia,
bruciando tonnellate di carburante per alzarsi di qualche centimetro e
vaporizzando l'acqua del mare perché fungesse da silenziatore di un
evento acustico che altrimenti li avrebbe scossi riducendoli a pezzi.
Poi fu come se avessero costruito delle scale di vapore e fumo e le
avessero salite, la loro velocità ora era visibile, pennacchi di fuoco
che sovrastavano i rotoli di nuvole che avevano creato. Andati e
spariti, come ogni altro lancio riuscito: rapido e vivido come un sogno,
e poi via, andati e spariti.

Ritardarono la partenza dell'ultimo razzo a causa di un sensore
difettoso, ma dieci minuti dopo anche quello fu lanciato con successo.
Sarebbe arrivato su Marte all'incirca un centinaio di anni dopo il resto
della flotta, ma questo era già stato preso in considerazione durante la
pianificazione e avrebbe potuto dimostrarsi positivo, un'iniezione di
tecnologia terrestre e di competenze molto tempo dopo che i libri di
carta e i lettori digitali dei primi coloni si sarebbero ridotti in
polvere.

ooo

Pochi istanti dopo la trasmissione si spostò sulla Guyana francese, in
particolare sul vecchio e ampliato Centro Nazionale di Studi Spaziali di
Kourou, dove uno dei razzi si era innalzato in cielo per una trentina di
metri e poi aveva perso la spinta ed era ricaduto sulla sua piattaforma
in un fungo di fuoco.

Dodici persone rimasero uccise, dieci a bordo dell'arca PNE e due a
terra, ma fu l'unica tragedia ingente dell'intera sequenza di lanci: una
gran fortuna, tutto considerato.

ooo

Ma quella non fu la fine dell'esercitazione. Verso mezzanotte - e fino a
quel momento mi parve l'indicatore più chiaro della disparità grottesca
tra il tempo terrestre e quello dello Spin - la civilizzazione umana su
Marte poteva aver fallito completamente o si stava svolgendo da più di
undicimila anni.

Quello era pressappoco il tempo intercorso tra la fine del Pleistocene e
il pomeriggio del giorno precedente, ed era passato mentre guidavo verso
il mio appartamento al ritorno dalla Perielio. Era del tutto possibile
che le dinastie marziane fossero fiorite e cadute mentre aspettavo che
il semaforo cambiasse colore. Pensai a tutte quelle vite - quelle vite
umane e reali, ognuna intrappolata nell'arco di un minuto o meno di
orologio - e mi sentii un po' stordito. Vertigine dello Spin. O qualcosa
di più profondo.

Quella notte furono lanciati una mezza dozzina di satelliti sonda
programmati per cercare segni di vita umana su Marte. I loro pacchetti
dati vennero quindi paracadutati sulla Terra e recuperati prima che
fosse mattina.

ooo

Vidi i risultati prima che fossero resi pubblici.

Ciò avvenne una settimana dopo i lanci Prometheus. Jason aveva prenotato
un appuntamento in infermeria per le 10.30, modificabile in caso di
notizie straordinarie dal JPL. Non cancellò l'appuntamento ma si
presentò in ritardo di un'ora con una busta manila in mano, chiaramente
ansioso di parlare di qualcosa che non aveva a che fare con la sua
terapia di medicinali. Ci spostammo in fretta in ambulatorio.

«Non so cosa dire alla stampa» disse. «Sono appena uscito da una
videoconferenza con il direttore dell'Agenzia Spaziale Europea e un
branco di burocrati cinesi. Stiamo cercando di mettere insieme una bozza
di dichiarazione congiunta per i capi di Stato, ma appena i russi sono
d'accordo con una frase i cinesi vogliono porre il veto e viceversa.»

«Una dichiarazione riguardo cosa, Jase?»

«I dati satellitari.»

«Hai i risultati?»

In effetti, erano attesi da un po' di tempo. Il JPL di solito era più
rapido nella condivisione delle fotografie. Ma da ciò che aveva detto
Jase, sospettavo che qualcuno stesse temporeggiando sui dati. Il che
significava che non erano ciò che si aspettavano. Cattive notizie,
forse.

«Guarda» disse Jason.

Aprì la busta manila e tirò fuori due foto telescopiche composite, una
sopra l'altra. Erano entrambe immagini di Marte prese dall'orbita della
Terra dopo i lanci Prometheus.

La prima foto era da arresto cardiaco. Non era così definita come
l'immagine incorniciata che avevo messo in sala d'attesa, giacché in
questa il pianeta era lontano dalla sua posizione più vicina possibile
alla Terra; la chiarezza che possedeva era una testimonianza della
moderna tecnologia delle immagini. Apparentemente non sembrava molto
diversa dalla foto incorniciata: riuscivo a distinguere abbastanza verde
da sapere che l'ecologia trapiantata era ancora intatta, ancora attiva.

«Guarda meglio» disse Jason.

Fece scorrere il dito lungo la linea sinuosa di una pianura fluviale.
C'erano zone verdi che avevano bordi netti e regolari. Più osservavo più
ne notavo.

«Agricoltura» disse Jase.

Trattenni il respiro e pensai a cosa volesse dire: ``Ora ci sono due
pianeti abitati nel sistema solare.'' Non si parlava più di ipotesi ma
di realtà. Erano posti in cui vivevano delle persone, in cui
\emph{vivevano} delle persone \emph{su Marte}.

Volevo stare lì a fissarla. Ma Jase infilò di nuovo la foto stampata
nella busta rivelando quella sottostante.

«La seconda foto,» disse, «è stata fatta ventiquattr'ore dopo.»

«Non capisco.»

«È stata scattata dalla stessa macchina fotografica dello stesso
satellite. Abbiamo immagini parallele a conferma dei risultati. Sembrava
un'imperfezione nel sistema di elaborazione d'immagini finché non
abbiamo aumentato il contrasto tanto da riuscire a registrare una
piccola luce stellare.»

Ma non c'era niente \emph{nella} foto. Qualche stella, un grosso nulla
centrale a forma di disco. «Cos'è?»

«Una membrana dello Spin» disse Jason. «Vista da fuori. Marte adesso ne
ha una sua.»
